% !TeX root = ../泛函分析笔记.tex
% 由主文档通过 \input 载入；请编译主文档。
\setcounter{chapter}{8} % 保留原书章号。
\chapter{弱紧性与自反空间}
\label{chap:fa-reflexive}

弱收敛的优势在于：某些不具有范数紧性的有界集合，在弱拓扑下恢复紧性。对偶空间中的一般结论是弱星紧性；要把它转化为原空间中的弱紧性，需要自反性。

\section{Banach--Alaoglu 定理}
先看一个不需要一般拓扑的模型：$\ell^2$ 中每个有界序列都有弱收敛子列。设 $\|x^{(n)}\|_2\leq M$，每个坐标序列有界。先选子列使第一坐标收敛，再从中选子列使第二坐标也收敛，依次进行；取这些嵌套子列的对角子列，使所有坐标同时收敛，记极限为 $x_j$。对每个有限 $N$，
\[
 \sum_{j=1}^N|x_j|^2=\lim_k\sum_{j=1}^N|x_j^{(n_k)}|^2\leq M^2.
\]
令 $N\to\infty$，可知 $x\in\ell^2$、$\|x\|_2\leq M$；再由命题\ref{prop:fa-coordinate-weak}得到弱收敛。一般弱紧性理论要把“逐坐标提取极限”的思想推广到由全部连续泛函提供的坐标，其中坐标集可能不可数。

\begin{theorem}\label{thm:fa-alaoglu}
任意赋范空间 $X$ 的对偶闭单位球 $B_{X^*}$ 在 $\sigma(X^*,X)$ 下紧。
\end{theorem}
\begin{proof}
将 $f\in B_{X^*}$ 视为坐标族 $(f(x))_{x\in X}$。每个坐标属于紧集 $D_x=\{z\in\mathbb K:|z|\leq\|x\|\}$，而乘积 $\prod_{x\in X}D_x$ 由 Tychonoff 定理紧。这个乘积中满足
\[
 a_{\alpha x+\beta y}=\alpha a_x+\beta a_y
 \quad(x,y\in X,\ \alpha,\beta\in\mathbb K)
\]
的坐标族构成闭集，因为每个等式都是有限坐标上的闭条件。它们恰好对应范数至多一的线性泛函。乘积拓扑在此闭集上的诱导拓扑正是弱星拓扑，所以 $B_{X^*}$ 紧。所用乘积紧性是附录\ref{app:fa-topology}中明确列出的拓扑输入。
\end{proof}
这一定理没有声称原空间 $X$ 的闭单位球弱紧，也没有在一般情形下声称每个有界泛函序列都有弱星收敛子列。

\section{可分空间中的弱星子列}
\begin{theorem}\label{thm:fa-weak-star-subsequence}
若 $X$ 可分，则 $X^*$ 中每个范数有界序列都有弱星收敛子列。
\end{theorem}
\begin{proof}
取可数稠密集 $\{x_j\}$。若 $\|f_n\|\leq M$，则对每个 $j$，标量列 $(f_n(x_j))$ 有界。逐次选子列并取对角子列，得到 $f_{n_k}(x_j)$ 对所有 $j$ 收敛。对任意 $x$，估计
\[
 |f_{n_k}(x)-f_{n_l}(x)|\leq2M\|x-x_j\|
       +|f_{n_k}(x_j)-f_{n_l}(x_j)|
\]
说明 $(f_{n_k}(x))$ 为 Cauchy 序列。令 $f(x)$ 为其极限，便得线性泛函及 $|f(x)|\leq M\|x\|$，所以 $f\in X^*$，子列弱星趋于 $f$。
\end{proof}
在 $B_{X^*}$ 上，还可用
\[
 d(f,g)=\sum_{j\geq1}2^{-j}\frac{|f(x_j)-g(x_j)|}{1+|f(x_j)-g(x_j)|}
\]
给弱星拓扑赋予度量。稠密性与统一范数界保证该度量收敛等价于逐点收敛。这里的有界集限制不可随意省略。

\section{自反空间与典范嵌入}
\begin{definition}
Banach 空间 $X$ 称为\emph{自反的（reflexive）}，若典范等距映射 $J:X\to X^{**}$，$(Jx)(f)=f(x)$，为满射。
\end{definition}
Hilbert 空间自反。实情形连续两次使用 Riesz 表示即可；复情形设 $R(y)=\langle\,\cdot\,,y\rangle$，对 $F\in H^{**}$，映射 $y\mapsto\overline{F(Ry)}$ 为连续线性泛函，故等于 $\langle y,z\rangle$。于是 $F(Ry)=\langle z,y\rangle=(Ry)(z)$，从而 $F=Jz$。

由定理\ref{thm:fa-sequence-dual}，$1<p<\infty$ 时 $\ell^p$ 自反；在本书采用的 $\sigma$ 有限测度空间上，$L^p$ 也由其对偶表示得出同一结论。另一方面，$c_0^{**}\cong\ell^\infty$，典范映射就是自然包含，因此 $c_0$ 不自反。自反性要求的是典范映射满射，而不只是抽象存在某个同构。

可以逐个追踪这个典范映射。一个 $x\in c_0$ 对 $a\in\ell^1=c_0^*$ 的作用是 $a\mapsto\sum_ja_jx_j$；而双对偶中还允许用任意有界数列 $b$ 定义 $a\mapsto\sum_ja_jb_j$。取 $b=\boldsymbol1$ 就得到双对偶中的元素，却没有 $c_0$ 中的原像。对 $1<p<\infty$，连续取两次共轭指数会回到 $p$，而上述配对把 $Jx$ 对应到原来的同一个数列 $x$，这才证明典范映射满射。只写两个空间“大小相同”并不足以证明自反性。

\section{自反性的弱紧性刻画}
\begin{lemma}[Goldstine 定理]\label{lem:fa-goldstine}
$J(B_X)$ 在 $B_{X^{**}}$ 中弱星稠密。
\end{lemma}
\begin{proof}
给定 $F\in B_{X^{**}}$ 及有限个 $f_1,\ldots,f_m\in X^*$，只需证明 $(F(f_j))_{j=1}^m$ 属于 $\{(f_j(x))_{j=1}^m:x\in B_X\}$ 的闭包。若不属于，在有限维实空间（复情形取实部）中分离点与闭凸集，得到 $f=\sum_jc_jf_j$ 满足
\[
 \operatorname{Re}F(f)>\sup_{x\in B_X}\operatorname{Re}f(x)=\|f\|.
\]
但 $\|F\|\leq1$ 保证 $\operatorname{Re}F(f)\leq\|f\|$，矛盾。有限坐标邻域恰好定义弱星拓扑，故引理成立。
\end{proof}
\begin{theorem}\label{thm:fa-reflexive-compact}
Banach 空间 $X$ 自反当且仅当 $B_X$ 弱紧。
\end{theorem}
\begin{proof}
自反时，$J$ 把 $B_X$ 双射到 $B_{X^{**}}$，并把弱拓扑对应到 $\sigma(X^{**},X^*)$；由 Banach--Alaoglu 定理得紧性。反之，若 $B_X$ 弱紧，则 $J(B_X)$ 在 Hausdorff 弱星拓扑中紧，因而闭；Goldstine 定理又保证它在 $B_{X^{**}}$ 中稠密，所以两者相等，缩放后得 $J(X)=X^{**}$。
\end{proof}
自反空间的闭线性子空间 $M$ 也自反：$B_M=M\cap B_X$ 为弱闭的弱紧集，而 Hahn--Banach 延拓保证 $M$ 的自身弱拓扑与诱导弱拓扑一致。闭子空间的商 $X/M$ 也自反：商映射弱连续，闭单位球的像弱紧；它包含商空间的开单位球，而后者在商闭单位球中弱稠密，故像包含整个商闭单位球，反向包含由商范数定义成立，再用定理\ref{thm:fa-reflexive-compact}。

\section{有界序列的弱收敛子列}
\begin{theorem}\label{thm:fa-reflexive-subsequence}
自反 Banach 空间中的每个有界序列都有弱收敛子列。
\end{theorem}
\begin{proof}
设 $(x_n)$ 有界，令 $M=\overline{\operatorname{span}\{x_n:n\geq1\}}$。$M$ 可分且自反。先说明 $M^*$ 可分：由上一节的度量和 Banach--Alaoglu 定理，$B_{M^*}$ 弱星紧且可度量化，故有可数弱星稠密子集。因 $M$ 自反，$M^*$ 上的弱星拓扑就是弱拓扑；该可数集的凸包弱稠密，定理\ref{thm:fa-convex-weak-closure}使其范数稠密。取有理系数的有限凸组合，便得到范数可数稠密集。

现在把 $Jx_n$ 看作 $(M^*)^*$ 中的有界序列，对可分空间 $M^*$ 使用定理\ref{thm:fa-weak-star-subsequence}，得到弱星收敛子列。极限由自反性等于某个 $Jx$，$x\in M$，故 $x_n$ 的相应子列在 $M$ 中弱收敛；将 $X^*$ 的泛函限制到 $M$，即知它也在 $X$ 中弱收敛。
\end{proof}
更一般的 Eberlein--Šmulian 定理说明：Banach 空间中，一个子集弱紧当且仅当其中每个序列都有弱收敛到该子集某点的子列。本书不证明这一更深结果；弱紧性的进一步理论可参阅\cite{teschl-fa}。上面的自反情形已足够支持后续直接法。

\section{例题与习题}
\begin{example}[最近点存在]
自反 Banach 空间中，任意非空闭凸集 $C$ 到给定点 $x$ 的距离都能达到。
\end{example}
\begin{solve}
取 $y_n\in C$ 使 $\|x-y_n\|\to d=\operatorname{dist}(x,C)$。此序列有界，故有弱收敛子列，记极限为 $y$。闭凸集弱闭，所以 $y\in C$；范数下半连续性给出 $\|x-y\|\leq d$，因而等号成立。唯一性通常需要严格凸性等额外条件。
\end{solve}

所谓严格凸，是指不同单位向量 $u,v$ 的中点满足 $\|(u+v)/2\|<1$。若两个不同点都是距离为 $d>0$ 的最近点，对它们相对于 $x$ 的误差向量归一化，严格凸性会使中点距离小于 $d$，矛盾；$d=0$ 时唯一性直接成立。Hilbert 范数严格凸，这与投影定理的唯一性相符。一般自反空间未必严格凸，例如 $\mathbb R^2$ 配最大值范数，到 $x$ 轴上点集的最近点可能形成一段线段。

\begin{example}[缺少自反性时的闭凸极小化问题]
在 $c_0$ 上令 $f(x)=\sum_{j\geq1}2^{-j}x_j$，并取闭仿射超平面 $C=\{x:f(x)=1\}$。则 $\inf_{x\in C}\|x\|_\infty=1$，但下确界不能达到。
\end{example}
\begin{solve}
因为 $\|f\|=\sum_j2^{-j}=1$，每个 $x\in C$ 都满足 $\|x\|_\infty\geq1$。令前 $N$ 个坐标均为 $(1-2^{-N})^{-1}$，其余为零，就得到 $x^{(N)}\in C$ 且范数趋于一。若某个 $x\in C$ 的范数为一，由 $1=\sum_j2^{-j}\operatorname{Re}x_j$ 及每个 $\operatorname{Re}x_j\leq1$，所有实部必须为一，结合 $|x_j|\leq1$ 得所有 $x_j=1$，不属于 $c_0$。因此闭凸约束和有界极小化序列本身还不够，还需要能保住极限的紧性条件。
\end{solve}
\begin{enumerate}
 \item 证明 $\ell^1$ 不自反。提示：若 $e_n$ 有弱收敛子列，坐标泛函迫使其极限为零，而求和泛函给出矛盾。
 \item\label{exer:fa-norm-attainment} 证明自反空间上的每个连续线性泛函都在闭单位球上达到其模的最大值。
 \item 证明严格凸 Banach 空间中，闭凸集的最近点一旦存在便唯一。
 \item 解释为何 $\ell^\infty$ 不自反：它含有闭子空间 $c_0$。
\end{enumerate}

